% Laget automatisk av elevark/lag_elevark.py fra sim/millikan.html. Endringer her blir overskrevet.
\documentclass[11pt]{article}
\usepackage{elevark}
\begin{document}
\elevarkhode{Fysikk 2 · Elektriske felt og krefter}{Millikans oljedråpeforsøk}{https://fysikklab.no/sim/millikan.html}{Simuleringen}

I 1909 sprayet Robert Millikan og Harvey Fletcher bittesmå oljedråper inn mellom to ladde plater. Noen dråper ble ladd av friksjonen i sprayen. Ved å justere spenningen kunne de få en dråpe til å sveve. Da er den elektriske kraften like stor som tyngden, og ladningen kan regnes ut. Alle ladningene de fant, var hele multipler av en minste ladning.

\section*{Oppgaver}

Start simuleringen med dråpe 1 og $U = 0\,\mathrm{V}$.

\begin{oppgaveliste}
\oppgave \fase{gjett}{Gjett}Dråpen faller når spenningen er null. Hva skjer når du skrur opp spenningen? Hvorfor må den øverste plata være positiv for at dette skal virke?
\svar
\oppgave \fase{forklar}{Forklar}Tegn kreftene på en dråpe som svever. Vis at ladningen er $q = mgd/U$.
\svar
\oppgave \fase{test}{Mål}Finn spenningen som får dråpe 1 til å sveve. Bruk knappene ±1 V til slutt. Regn ut massen og ladningen til dråpen, og skriv tallene inn i tabellen.
\svar
\oppgave \fase{test}{Mål}Gjør det samme for alle de åtte dråpene. Hva legger du merke til med ladningene?
\svar
\oppgave \fase{forklar}{Forklar}Finn den største ladningen som alle ladningene er et helt antall ganger. Hvor mange elementærladninger har hver dråpe? Kryss av for linjene for å sjekke.
\svar
\oppgave \fase{forklar}{Vurder}Millikan målte ikke radien direkte. Hvorfor er det vanskelig å måle radien til en dråpe på under 1~µm med et mikroskop? Hvor mye feil gir 5~\% feil i radien i ladningen?
\svar
\oppgave \fase{utfordring}{Utfordring}Når $U = 0$, faller dråpen med terminalfart. Luftmotstanden er da $6πηrv$, der η = $1{,}81 \cdot 10^{-5}\,\mathrm{Pa}$ $s$. Velg ekte tid, mål hvor lang tid dråpe 1 bruker på å falle 0,1~mm, og bruk det til å regne ut radien. Stemmer det med radien i simuleringen?
\svar
\end{oppgaveliste}

\end{document}
